% GENERATED FILE. Edit canonical lessons, then run npm run generate:books.
% canonical-source-hash: fnv1a64:395f5e511f2065b0
% canonical-lessons: GE-C470-geduld, GE-C470-gelassen, GE-C470-beruhigen, GE-C470-aufregung, GE-C470-zappeln

\chapter{Geduld und Gelassenheit}
\label{ch:ge-geduld-und-gelassenheit}
\hlchaptermodality{470}

\begin{chapteropening}
\textbf{By the end of this chapter:} \emph{I can talk about patience and staying calm.}

Complete the last lesson of chapter 470: I can talk about patience and staying calm.
\end{chapteropening}

\section[die Geduld]{die Geduld — patience}

\label{lesson:GE-C470-geduld}



\begin{quote}
Before the new one: say the German for to breathe a sigh of relief, then the German for to sigh.
\end{quote}

\subsection*{You'll want to know: die Geduld}
\textbf{die Geduld} — \textquotedblleft{}patience\textquotedblright{}.

\begin{quote}
Mit kleinen Kindern braucht man viel Geduld.
\end{quote}

\begin{quote}
\emph{With small children you need a lot of patience.}
\end{quote}

No plural. \emph{Die Geduld verlieren} — to lose patience.

\begin{cousinweb}
\emph{Geduld} is inherited from Middle High German \emph{gedult} and Old High German \emph{gidult} (EtymDB), usually traced to Proto-Germanic \emph{gaþuldiz}; \emph{dulden} \textquotedblleft{}to endure\textquotedblright{} is related.
\end{cousinweb}

\begin{grammarlens}[title={What you've built}]
One more word for telling, explaining and describing at B1.
\end{grammarlens}

\subsection*{Your turn}
\begin{itemize}
\raggedright
  \item \emph{Say it:} \emph{die Geduld}
  \item \emph{Say it:} \emph{die Geduld}, once more
  \item \emph{Say it:} a sentence of your own with \emph{die Geduld}
  \item \emph{Recall:} say the German for to breathe a sigh of relief, then the German for to sigh, then say \emph{die Geduld} again
\end{itemize}

\begin{tcolorbox}[breakable,colback=teal!4,colframe=teal!35!black,title={Before you move on}]
What does die Geduld mean? (\textquotedblleft{}Patience\textquotedblright{}.) Say it once more.
\end{tcolorbox}
\section[gelassen]{gelassen — calm, composed}

\label{lesson:GE-C470-gelassen}



\begin{quote}
Before the new one: say the German for to sigh, then the German for patience.
\end{quote}

\subsection*{You'll want to know: gelassen}
\textbf{gelassen} — \textquotedblleft{}calm, composed\textquotedblright{}.

\begin{quote}
Trotz des Stresses blieb mein Chef ganz gelassen.
\end{quote}

\begin{quote}
\emph{Despite the stress, my boss stayed quite calm.}
\end{quote}

Also the past participle of \emph{lassen}: \emph{Ich habe den Schirm im Bus gelassen.}

\begin{cousinweb}
\emph{gelassen} is the old past participle of \emph{lassen} \textquotedblleft{}to let, to let go\textquotedblright{}, which is inherited from Old High German \emph{lāzan}, from Proto-Germanic \emph{lētaną} (EtymDB): someone who has let things go.
\end{cousinweb}

\begin{grammarlens}[title={What you've built}]
One more word for telling, explaining and describing at B1.
\end{grammarlens}

\subsection*{Your turn}
\begin{itemize}
\raggedright
  \item \emph{Say it:} \emph{gelassen}
  \item \emph{Say it:} \emph{gelassen}, once more
  \item \emph{Say it:} a sentence of your own with \emph{gelassen}
  \item \emph{Recall:} say the German for to sigh, then the German for patience, then say \emph{gelassen} again
\end{itemize}

\begin{tcolorbox}[breakable,colback=teal!4,colframe=teal!35!black,title={Before you move on}]
What does gelassen mean? (\textquotedblleft{}Calm, composed\textquotedblright{}.) Say it once more.
\end{tcolorbox}
\section[beruhigen]{beruhigen — to calm (down)}

\label{lesson:GE-C470-beruhigen}



\begin{quote}
Before the new one: say the German for patience, then the German for calm.
\end{quote}

\subsection*{You'll want to know: beruhigen}
\textbf{beruhigen} — \textquotedblleft{}to calm (down)\textquotedblright{}.

\begin{quote}
Die Mutter beruhigt das weinende Baby.
\end{quote}

\begin{quote}
\emph{The mother calms the crying baby.}
\end{quote}

\emph{Sich beruhigen} — to calm down: \emph{Er hat sich schnell beruhigt.}

\begin{cousinweb}
\emph{beruhigen} is built on \emph{be-} and \emph{ruhig} \textquotedblleft{}calm\textquotedblright{}, from \emph{Ruhe} \textquotedblleft{}rest\textquotedblright{}, which is inherited from Old High German \emph{ruowa}, from Proto-Germanic \emph{rōwō} (EtymDB).
\end{cousinweb}

\begin{grammarlens}[title={What you've built}]
One more word for telling, explaining and describing at B1.
\end{grammarlens}

\subsection*{Your turn}
\begin{itemize}
\raggedright
  \item \emph{Say it:} \emph{beruhigen}
  \item \emph{Say it:} \emph{beruhigen}, once more
  \item \emph{Say it:} a sentence of your own with \emph{beruhigen}
  \item \emph{Recall:} say the German for patience, then the German for calm, then say \emph{beruhigen} again
\end{itemize}

\begin{tcolorbox}[breakable,colback=teal!4,colframe=teal!35!black,title={Before you move on}]
What does beruhigen mean? (\textquotedblleft{}To calm (down)\textquotedblright{}.) Say it once more.
\end{tcolorbox}
\section[die Aufregung]{die Aufregung — excitement; agitation}

\label{lesson:GE-C470-aufregung}



\begin{quote}
Before the new one: say the German for calm, then the German for to calm.
\end{quote}

\subsection*{You'll want to know: die Aufregung}
\textbf{die Aufregung} — \textquotedblleft{}excitement; agitation\textquotedblright{}.

\begin{quote}
Vor lauter Aufregung konnte sie nicht schlafen.
\end{quote}

\begin{quote}
\emph{She was so excited she could not sleep.}
\end{quote}

Built on \emph{sich aufregen}. Plural, rare: \emph{die Aufregungen}.

\begin{cousinweb}
\emph{Aufregung} is built on \emph{aufregen} with \emph{-ung}; \emph{regen} \textquotedblleft{}to stir\textquotedblright{} goes back to Middle High German \emph{regen} \textquotedblleft{}to excite, to move\textquotedblright{} (EtymDB).
\end{cousinweb}

\begin{grammarlens}[title={What you've built}]
One more word for telling, explaining and describing at B1.
\end{grammarlens}

\subsection*{Your turn}
\begin{itemize}
\raggedright
  \item \emph{Say it:} \emph{die Aufregung}
  \item \emph{Say it:} \emph{die Aufregung}, once more
  \item \emph{Say it:} a sentence of your own with \emph{die Aufregung}
  \item \emph{Recall:} say the German for calm, then the German for to calm, then say \emph{die Aufregung} again
\end{itemize}

\begin{tcolorbox}[breakable,colback=teal!4,colframe=teal!35!black,title={Before you move on}]
What does die Aufregung mean? (\textquotedblleft{}Excitement; agitation\textquotedblright{}.) Say it once more.
\end{tcolorbox}
\section[zappeln]{zappeln — to fidget, to wriggle}

\label{lesson:GE-C470-zappeln}



\begin{quote}
Before the new one: say the German for to calm, then the German for excitement.
\end{quote}

\subsection*{You'll want to know: zappeln}
\textbf{zappeln} — \textquotedblleft{}to fidget, to wriggle\textquotedblright{}.

\begin{quote}
Der kleine Junge zappelt auf seinem Stuhl.
\end{quote}

\begin{quote}
\emph{The little boy is fidgeting on his chair.}
\end{quote}

\emph{Jemanden zappeln lassen} — to keep someone on tenterhooks.

\begin{cousinweb}
\emph{zappeln} is usually traced to Middle High German \emph{zabelen} \textquotedblleft{}to wriggle\textquotedblright{}; EtymDB lists the German word without a source, and its older origin is not settled.
\end{cousinweb}

\begin{grammarlens}[title={What you've built}]
One more word for telling, explaining and describing at B1.
\end{grammarlens}

\subsection*{Your turn}
\begin{itemize}
\raggedright
  \item \emph{Say it:} \emph{zappeln}
  \item \emph{Say it:} \emph{zappeln}, once more
  \item \emph{Say it:} a sentence of your own with \emph{zappeln}
  \item \emph{Recall:} say the German for to calm, then the German for excitement, then say \emph{zappeln} again
\end{itemize}

\begin{tcolorbox}[breakable,colback=teal!4,colframe=teal!35!black,title={Before you move on}]
What does zappeln mean? (\textquotedblleft{}To fidget, to wriggle\textquotedblright{}.) Say it once more.
\end{tcolorbox}
